\documentclass[a4paper,11pt,twopage,numbers=noenddot]{scrbook}

\usepackage[top=2cm,lmargin=1in,rmargin=1in,bottom=3cm,hmarginratio=1:1]{geometry}
\usepackage[utf8]{inputenc}
\usepackage[english]{babel}
\usepackage{amssymb}
\usepackage{amsthm}
\usepackage{thmtools}
\usepackage{fancyvrb}
\usepackage{mathtools}
\usepackage{amsmath}
\usepackage{ifthen}
\usepackage{xspace}
\usepackage{import}
\usepackage{calc}
\usepackage{float}
\usepackage{listings}
\usepackage[pdftex,
  pdfauthor={Michael Banken},
  pdftitle={Simulation of countermachines},
] {hyperref}
\usepackage{makeidx}
\usepackage{graphicx}
%\usepackage[right]{showlabels}
%\usepackage[justific=raggedright,totoc]{idxlayout}

\graphicspath{{img/}}

\addto\extrasenglish{%
  \renewcommand{\chapterautorefname}{Section}
  \renewcommand{\sectionautorefname}{Section}
  \renewcommand{\subsectionautorefname}{Subsection}
  }

\newcommand\chap[1]{%
  \chapter*{#1}%
  \chaptermark{#1}%
  \addcontentsline{toc}{chapter}{#1}}

\declaretheorem[name=Definition,style=definition,numberwithin=chapter]{definition}
\declaretheorem[name=Example,style=definition,sibling=definition]{example}
\declaretheorem[style=definition,numbered=no]{exercise}
\declaretheorem[name=Remark,style=definition,sibling=definition]{remark}
\declaretheorem[name=Assumption,style=definition,sibling=definition]{assumption}
\declaretheorem[name=Observation,style=definition,sibling=definition]{observation}
\declaretheorem[name=Theorem,sibling=definition]{theorem}
\declaretheorem[sibling=definition]{corollary}
\declaretheorem[name=Fact,sibling=definition]{fact}
\declaretheorem[sibling=definition]{lemma}
\declaretheorem[sibling=lemma]{proposition}
\declaretheorem[name=Instruction set,style=definition,numberwithin=chapter]{instructionset}

%%%%%%%%%%%%%%%%%%%%%%%%%%%%%%%%%%%%%%%%%%%%%%%%%%%%%%%%%%%%%%%%%%%%%%%%%%%
%                                                                         %
%                   Spacing settings                                      %
%                                                                         %
%%%%%%%%%%%%%%%%%%%%%%%%%%%%%%%%%%%%%%%%%%%%%%%%%%%%%%%%%%%%%%%%%%%%%%%%%%%
\setlength{\parindent}{0pt}
\setlength{\parskip}{6pt}
\setlength{\marginparsep}{0cm}

% kill those ugly red rectangles around links
\hypersetup{
  colorlinks,%
  citecolor=black,%
  filecolor=black,%
  linkcolor=black,%
  urlcolor=black,
}

\usepackage{newunicodechar}
\newunicodechar{∀}{\ensuremath{\forall}}
\newunicodechar{∃}{\ensuremath{\exists}}
\newunicodechar{×}{\ensuremath{\times}}
\newunicodechar{ø}{\ensuremath{\emptyset}}
\newunicodechar{≤}{\ensuremath{\le}}
\newunicodechar{∈}{\ensuremath{\in}}
\newunicodechar{→}{\ensuremath{\to}}
\newunicodechar{⊆}{\ensuremath{\subseteq}}
\newunicodechar{⊗}{\ensuremath{\otimes}}
\newunicodechar{∧}{\ensuremath{\wedge}}


\begin{document}
\pagestyle{plain}
\begin{titlepage}
\newcommand{\drop}{0.07\textheight}
\begin{center}
\begingroup%
\vfill
{\LARGE\textsc{%
Friedrich-Alexander-Universität\\[2mm]
Erlangen-Nürnberg%
}}\\[\drop]
{%
% trim= left bottom right top
\includegraphics[height=1.8cm,trim=0cm 0mm 0 0mm]{img/tcs}\\
\textsc{\large Chair for Computer Science 8}\\
\textsc{\large Theoretical Computer Science}}%\\[\drop]
\vfill
\rule{\textwidth}{1pt}\par
\vspace{0.5\baselineskip}
{% maybe \itshape?
\Huge\bfseries Simulation of countermachines
\\[1cm]
\large\bfseries
A simulator for the Minsky machine
\\[1cm]
\textbf{\large Bachelor Thesis in Computer Science}
}\\[0.5\baselineskip]
\rule{\textwidth}{1pt}\par
\vfill
{\Large{Michael Banken}\\
%{\large\normalsize{michael.banken@fau.de}}
}
\vfill
\large Advisors:
\vfill
\begin{tabular}{ccc}
    \large
    Tadeusz Litak
\end{tabular}
\vfill
% trim= left bottom right top
\includegraphics[height=1.8cm,trim=0cm -5mm 0 0mm]{img/fau}
\vfill
{\large Erlangen, \today}
\endgroup
\end{center}
\end{titlepage}

% vim: tw=80 spell spelllang=en nocul
%
\tableofcontents
\chapter{Introduction}
\section{Historical context}
\section{Motivation}

\chapter{Theoretical background}
\section{Counter machines}

\section{Minsky machine}
One counter machine that is of particular interest for this work is the Minsky machine, named after Marvin L. Minsky who proposed it back in 1961. Minsky originally described these machines in a paper discussing the recursive unsolvability of \cite{10.2307/1970290}. In this paper he defines them as a restricted form of the 2-band Turing machine, as follows below.
\begin{definition}{Two tape non-writing Turing machine:}
The two tape non-writing Turing machine consists of
\begin{itemize}
\item a finite state machine
\item 2 semi infinite bands
\end{itemize}
Both bands are initialized with a symbol marking their respective ends and are otherwise filled blanks.
The machine can neither write, nor delete the symbols on either band.
\end{definition}
This means the state of the machine can be effectively described by a tuple $(p_0, p_1, s)$ with $p_0$, $p_1$ being the bands respective head positions and $s$ being the state of the finite state machine. See \autoref{fig2bandstate} for a graphical example of a 2-band Turing machine. The example machine will move $p_0$ to the right and $p_1$ to the left until $p_1$ is 0.
\begin{figure}[H]
	\def\svgscale{0.3}
	\input{img/twobandturingmachine.pdf_tex}
	\caption{visualisation of 2-band TM}
	\label{fig2bandstate}
\end{figure}
Minsky showed that any given Turing machine T can be represented by a two tape non-writing Turing machine T*\cite{10.2307/1970290}. He does this by first converting T to the equivalent 2-symbol Turing machine T' which can be obtained as described by Shannon\cite{Shannon1971-SHAAUT-2}. Then the resulting machine T' is converted to a two tape non-writing Turing machine. This thesis will not discuss the exact nature of this proof, but it can be found in the original paper. However we will utilize a similar idea  when we discuss how to write programs for the Minsky machine.
It is immediately obvious that with the machine not being able to write or write on the tape, the machine can only move either of its heads left or right and check if it has reached the end of a tape. This means a more natural approach to viewing this machine would be to represent the tapes as counters $x_1$ and $x_2$. Accordingly the state machine could be represented as a program counter $I_p$ and a program $p$ represented by a vector of subroutines formed from the following instruction set:
\begin{instructionset}
\label{minsky4instruction}
\begin{itemize}
\hfill
\item{$INC:$} $x_y = x_y + 1$; | $ y\in{0,1}$\hfill\break -- i.e. increment counter $x_y$
\item{$DEC:$} if $x_y > 0$ then: $x_y = x_y - 1$; | $ y\in{0,1}$\hfill\break -- i.e. decrement counter $x_y$ if possible
\item{$JMP:$} $I_p = i$; | $ 0 \leq i \leq sizeof(p)$\hfill\break -- i.e. jump to $i$
\item{$JIZ:$} if $x_y=0$ then: $ I_p = j$, else: $I_p = i$; | $y\in{0,1} \land 0 \leq i,j \leq sizeof(p)$\hfill\break -- i.e. branch to $j$ if counter $x_y$ is zero otherwise jump to $i$
\end{itemize}
\end{instructionset}
This rather simple instruction set can be simplified further to just two basic instructions. One such instruction set will be given below and we will focus on a Minsky machine defined with just these two instructions for the remainder of this thesis.

\subsection{Syntax and semantics}

For the purpose of this thesis we will define the Minsky machine as follows:
\begin{definition}
A Minsky machine is a tuple $(x_0, x_1, p, I_p)$, where $x_0$ and $x_1$ are counters, $I_p$ is the program counter and $p$ is a program consisting of the instructions found in the following \autoref{minsky2instruction}:
\begin{instructionset}
\label{minsky2instruction}
\hfill
\begin{itemize}
\item{$INC:$} $x_y = x_y + 1; I_p = i$; | $ y\in{0,1} \land 0 \leq i \leq sizeof(p)$\hfill\break -- i.e. increment counter $x_y$ and jump to $i$
\item{$DEC:$} if $x_y > 0$ then: $x_y = x_y - 1; I_p = i$, else: $I_p = j$; | $ y\in{0,1} \land 0 \leq i,j \leq sizeof(p)$\hfill\break -- i.e. decrement counter $x_y$ and jump to $i$ if possible, branch to $j$ if not
\end{itemize}
\end{instructionset}
\end{definition}
Before we can define the semantics for computing on our Minsky machine, we will first define denotational semantics for \autoref{minsky2instruction}
\begin{definition}

\end{definition}
\begin{lemma}
\label{lemma2btmtom4m}
Any 2-band Turing machine can be represented by an equivalent machine $(x_0, x_1, p, I_p)$ where $x_0$ and $x_1$ are counters, $I_p$ is the program counter and $p$ is a program consisting of the instructions found in \autoref{minsky4instruction}.
\begin{proof}
%TODO 
insert proof here
\end{proof}
\end{lemma}
\begin{theorem}
Any 2-band Turing machine can be represented by an equivalent Minsky machine.
\begin{proof}
In order to prove the existence of an equivalent Minsky machine we will show that any program from \autoref{minsky4instruction} has an equivalent program from \autoref{minsky2instruction} and vice versa. After that we can use \autoref{lemma2btmtom4m} to show that that for every 2-band Turing machine there is an equivalent Minsky machine.

%TODO 
insert proof here
\end{proof}
\end{theorem}
From this theorem immediately follows, due to the Turing-completeness of the 2-band Turing machine, that the Minsky machine is also Turing complete.
\begin{corollary}
The Minsky machine is Turing complete.
\end{corollary}

\subsection{Subroutines for the Minsky machine}
In order to make writing programs for the Minsky machine easier it is helpful to introduce the notion of subroutines.
For this we will introduce the notion of subroutines in a manner similar to how it was suggested by Shepherdson and Sturgis in 1963\cite{Shepherdson:1963:CRF:321160.321170}.
First we need to extend our instruction set by a call instruction:
\begin{instructionset}
\label{minskycallinstruction}
\hfill
\begin{itemize}
\item{$INC:$} $x_y = x_y + 1; I_p = i$; | $ y\in{0,1} \land 0 \leq i \leq sizeof(p)$\hfill\break -- i.e. increment counter $x_y$ and jump to $i$
\item{$DEC:$} if $x_y > 0$ then: $x_y = x_y - 1; I_p = i$, else: $I_p = j$; | $ y\in{0,1} \land 0 \leq i,j \leq sizeof(p)$\hfill\break -- i.e. decrement counter $x_y$ and jump to $i$ if possible, branch to $j$ if not
\item{$CALL:$} $subroutine$ $(registers)$ $[exits]$ | $registers$ and $exits$ are lists of positive integers and $subroutine$ is a positive integer
\end{itemize}
\end{instructionset}

Since the desired behaviour of this new instruction is clearly dependent on other parts of the machine we will first extend the definition of the Minsky machine\ref{smm} to account for the existence of subroutines and then give denotational semantics for the new instruction set in the context of one such machine.
\begin{definition}
\label{smm}
A subroutine Minsky machine is a tuple $(x_0, x_1, p', I_p, s)$, where $x_0$ and $x_1$ are counters, $I_p$ is the program counter and $p'$ is a program consisting of the instructions found in \autoref{minskycallinstruction} and $s$ is a list of subroutine programs also consisting of instructions from \autoref{minskycallinstruction}.
\end{definition}



\subsection{Writing programs for the Minsky machine}
While we now know that the Minsky machine is Turing complete, in practice writing programs can be challenging.
The main problem seems obvious as there are only two registers to work with and a very limited instruction set to do so.
Fortunately Minsky offered a solution by correctly encoding the input and output for the Minsky machine\cite{Minsky:1967:CFI:1095587}.



\chapter{The simulator}
\label{chap:simulator}
\section{Syntax and semantics}
\section{Usage of the simulator}
\subsection{None-graphical mode}
\subsection{Graphical user interface}
\subsection{Requirements}
\section{Implementation details}
\subsection{Representing the program}
In order to speed up the simulation the program is internally transformed into object form before execution. The parsing process will be explained in more detail below.\\
Subroutine calls are not represented in this form, since they are fully resolved by the parser. This may however not be an ideal solution due to problems we will discuss in \autoref{limit:program}, instead it may be a good idea to represent them as objects in future implementations, as will be discussed in \autoref{future:call}

\subsection{Parsing the program}
The parser sequentially reads the program and whenever it successfully recognizes a command it creates an equivalent Simulator\_command object, which is then added to the vector representing the program of the current subroutine.\\
An exception to this is the call of a subroutine, in which case the parser inserts the entire program of that subroutine, making all necessary substitutions as specified in the call. This copy and substitution process is carried out by an object implementing the visitor pattern and needs to be implemented for each specific specialization of the Simulator\_command class, which is something to keep in mind when adding additional instructions to the simulator.

\subsection{Representing the machine's state}
As we discussed in the previous chapter the state of a Minsky machine can be thought of as a triple $(x_0, x_1, I_p)$ with $x_0$, $x_1$ being the counters and $I_p$ being the program counter. Since these are merely three positive integers, we can easily model them as such.

Since $I_p$ will be used to reference structures lying in the memory of the host system, it is obviously best described by the size\_t type, on most 64-bit platforms this will be an unsigned 64-bit integer. By definition any value that cannot by represented as a size\_t is indexing a program that is too large for the host system, so the size limitations should not be a problem.

The counters are a slightly more complicated matter. The counters of the theoretical machine are by definition potentially infinitely large and many of the machine's theoretic properties rely strongly on this fact. Obviously this behaviour cannot be reproduced in a real simulator. It seems immediately obvious that big integers would allow to come reasonably close, so they may seem a natural way to implement the counters, however this is not the case with this simulator, instead counters are implemented as simple variables of type unsigned long long int. At first glance this may seem rather limited, but as we will discuss in more detail in \autoref{limit:registers}.
\subsection{Executing programs}
 
\subsection{Differences between the theoretical model and the simulated machine}
\subsection{Graphical Library}

\chapter{Limitations of simulating counter machines}
\label{chap:limitations}
\section{Memory}
\subsection{Registers}
\label{limit:registers}
\subsection{Program}
\label{limit:program}
Due to the nature of the internal representation of programs we also need to take into account the potential blow-up of storing these programs.
While the internal simulator commands themselves are all of a constant size, the nature of our subroutine substitution means that a single line of code of the original program can be expanded into an entire subroutine of internal commands. If that subroutine also contains calls to other subroutines it is easy to see how the internal representation of a program in execution could demand the use of vastly more memory than the program in its textual form did.
In order to simplify the calculations we will not focus on the actual text characters, but rather the number commands the program contains, neglecting the length of these commands and other syntax elements such as labels, routine headers and config commands.
In order to get a rough idea of how large this blow-up could get at its worst, let us make some very pessimistic assumptions.
Let n be the number of commands in the program, including all subroutines.
This means that there can be only n non-empty subroutines.
Since these routines can only call subroutines defined before them, the n-th routine can only call the n-1 subroutines before it.
From this follows, considering that any subroutine cannot contain more than n instructions, that the n-th routine can at worst make n calls to the largest subroutine before it.
Through simple induction over n this means that the n-th routine can at most expand to $n^{n}$ instruction. If this routine is the main routine it means that the entire program expands from n instructions in the source code to $n^{n}$ instructions in the simulators internal representation. This is of course less than ideal.

Fortunately these assumptions cannot possibly be true all at once, but the example code \ref{exponentialcode} below will show that at least an exponential blow-up is entirely possible.
\lstset{escapeinside={@(}{)@}}
\begin{lstlisting}[frame=single, caption=code resulting in exponential blow-up]
@(\label{exponentialcode})@
def subroutine1 [next]
	add 1
	add 1 next

def subroutine2 [next]
	call subroutine1 s
s:	call subroutine1 next

def subroutine3 [next]
	call subroutine2 s
s:	call subroutine2 next

...

def subroutine@($\frac{n}{2}-1$)@ [next]
	call subroutine@($\frac{n}{2}-2$)@ s
s:	call subroutine@($\frac{n}{2}-2$)@ next

main
	call subroutine@($\frac{n}{2}-1$)@ s
s:	call subroutine@($\frac{n}{2}-1$)@ 0
\end{lstlisting}
It is immediately obvious that this program consists of n instructions that will expand into $2^{\frac{n}{2}}$ instructions when parsed into the internal representation. It is easy to see that one can construct a similar program for any other natural number.
This means that the maximum blowup lies somewhere between $2^{\frac{n}{2}}$ and $n^{n}$.
\section{Runtime}
\subsection{Simulator performance}
\subsection{Computational complexities of Minsky programs}

\chapter{Future Work}
\label{chap:future}
\section{Adding additional instruction sets}
\section{Improving memory efficiency through more efficient representation of programs}

\chapter{Conclusions}
\label{chap:conclusions}

\bibliography{bib}{}
\bibliographystyle{plain}
\end{document}

% vim: tw=80 nospell spelllang=en nocul
